\chapter{Foundations \& Related Work}
\label{ch:foundations}

This chapter has two jobs. Sections~\ref{sec:fc-parcellation}--\ref{sec:survival}
establish the technical vocabulary the rest of the thesis uses. Section~\ref{sec:gaps}
does the harder and more important work: it states, as precisely as the current
literature allows, \emph{what is actually unsettled} in learning from brain functional
connectomes --- and where the experiments in this thesis land relative to those open
questions. That section is deliberately the longest in the chapter, because the central
empirical results of this work (Chapters~\ref{ch:sota-repair} and~\ref{ch:results}) are
only legible as contributions once the reader knows which of the field's assumptions
they contradict.

%==============================================================================
\section{Brain Connectivity \& Parcellation Schemes}
\label{sec:fc-parcellation}

Resting-state functional MRI measures the blood-oxygen-level-dependent (BOLD) signal, a
haemodynamic proxy for neural activity, while the subject performs no task. The standard
reduction from a 4-D BOLD volume to an analysable object proceeds in two steps. First, a
\emph{parcellation} partitions the cortex (and optionally subcortex) into $R$ regions of
interest (ROIs), collapsing $\sim 10^5$ voxels into $R$ representative time series. Second,
a similarity measure applied pairwise across those $R$ time series yields an $R \times R$
\emph{functional connectivity} (FC) matrix. Pearson correlation followed by the Fisher
$z$-transform is the near-universal default, and is what this work uses.

The choice of $R$ and of parcellation atlas is not a neutral preprocessing detail. The
Schaefer atlas family provides gradient-weighted, resting-state-derived parcellations at
multiple granularities (100, 200, 400, \ldots\ parcels), each parcel additionally assigned
to one of the canonical Yeo functional networks --- default mode (DMN), limbic, dorsal
attention, salience, and so on. Anatomically defined alternatives such as Harvard--Oxford
and AAL remain common in the clinical literature. This thesis uses \textbf{Schaefer-200}
throughout, which fixes $R = 200$ and therefore fixes the dimensionality of every
downstream representation: each subject-visit is a $200 \times 200$ symmetric matrix, and
each ROI's feature vector is its own row of that matrix --- its correlation profile
against all other regions.

That last identification matters more than it first appears, and Section~\ref{sec:gaps}
returns to it: in the graph formulation used here and in most of the literature, the node
features \emph{are} the connectivity data. The graph's edges are then derived from the very
same matrix. Node features and topology are therefore not two independent information
sources; they are two views of one object.

The DMN is of particular interest in Alzheimer's disease. It is the network whose
posterior hubs (posterior cingulate, precuneus) overlap most closely with the earliest
sites of amyloid deposition, and DMN dysconnectivity is among the most replicated
functional findings across the AD continuum. This motivates the region-restricted
variants of the pipeline discussed in Chapter~\ref{ch:conclusion} as future work.

%==============================================================================
\section{Graph Representation Learning in Neuroimaging}
\label{sec:gnn}

Treating an FC matrix as a graph $G = (V, E)$ with $|V| = R$ nodes invites the machinery
of graph neural networks (GNNs). The dominant paradigm is \emph{message passing}: each
node's representation is iteratively updated by aggregating transformed representations of
its neighbours,
\[
  h_v^{(\ell+1)} \;=\; \phi\!\left(h_v^{(\ell)},\ \bigoplus_{u \in \mathcal{N}(v)}
    \psi\!\left(h_v^{(\ell)}, h_u^{(\ell)}, e_{uv}\right)\right),
\]
where $\bigoplus$ is a permutation-invariant aggregator. Graph convolutional networks
(GCN) instantiate this with a degree-normalised mean; graph attention networks (GAT, and
its corrected successor GATv2) instead learn the aggregation weights, so that a node can
attend more strongly to some neighbours than others. GATv2 with FiLM-style conditioning is
the encoder architecture used in this thesis.

Because labelled clinical cohorts are small, a second strand of work borrows the
self-supervised pretraining recipe from language and vision: train an encoder on a
\emph{pretext} task that needs no labels, then reuse its representations downstream. For
graphs, the canonical pretext task is reconstruction. A graph autoencoder (GAE) encodes
each node to a latent $z_v$ and is trained so that a decoder can recover the adjacency
structure; the variational form (VGAE) additionally regularises the latent distribution
toward a prior via a Kullback--Leibler term. The premise --- that an encoder which can
rebuild a connectome must have learned something transferable about connectomes --- is
intuitive and near-universally assumed. Section~\ref{sec:gaps} explains why it is also
one of the field's least tested assumptions, and Chapter~\ref{ch:results} tests it
directly.

Benchmarking efforts such as BrainGB have attempted to standardise comparison across
brain-GNN architectures. Their existence is itself informative: the field reached a state
where dozens of architectures reported wins on incomparable pipelines before anyone
established what a common evaluation would even look like.

%==============================================================================
\section{Longitudinal \& Spatiotemporal Modelling}
\label{sec:longitudinal}

A single connectome is a snapshot. Disease is a trajectory. Modelling the AD continuum
therefore requires consuming a \emph{sequence} of per-visit connectomes, and the
literature offers two broad families.

\paragraph{Recurrent approaches.} A recurrent network over per-visit embeddings is the
parameter-lean option. The complication in clinical data is that visits are not evenly
spaced: a subject may be scanned at months $0$, $14$, and $39$, while another is scanned
at $0$, $12$, and $24$. A vanilla LSTM sees only ordinal positions $1, 2, 3$ and is blind
to that difference. Time-aware variants address this explicitly --- T-LSTM decomposes the
cell memory into short- and long-term components and applies an elapsed-time decay to the
short-term part, while ATTAIN combines time-awareness with attention over prior visits.
This thesis takes the simpler and more transparent route of conditioning the gate updates
on the inter-visit interval $\Delta t$ directly (Section~\ref{sec:gelstm-method}); the
design intent is identical.

\paragraph{Spatiotemporal graph transformers.} The higher-capacity option applies
self-attention across both space (ROIs) and time (visits). Brain-TokenGT, the published
baseline audited in Chapter~\ref{ch:sota-repair}, tokenises node and spatio-temporal edge
embeddings and feeds them to a transformer encoder augmented with trainable type
identifiers and non-trainable node identifiers. Its temporal component derives from
EvolveGCN-H, in which a gated recurrent unit evolves the \emph{weight matrix} of a graph
convolution rather than a hidden state --- an elegant idea whose optimisation behaviour on
small cohorts turns out to be the subject of Chapter~\ref{ch:sota-repair}.

\paragraph{Foundation models.} The current frontier scales pretraining rather than
architecture. Brain-JEPA applies a joint-embedding predictive architecture to fMRI with
brain-gradient positional encoding and spatiotemporal masking, pretrained on UK Biobank
and reporting state-of-the-art transfer across demographic, cognitive, and diagnostic
tasks; brain-graph foundation models pursue the same goal across atlases and disorders.
These models are trained on $10^4$--$10^5$ subjects. They are the correct point of
comparison for what large-scale pretraining buys, and the wrong point of comparison for
what is achievable on a single-site cohort of $N < 150$ --- a distinction
Section~\ref{sec:gaps} argues the field routinely blurs.

%==============================================================================
\section{Time-to-Event \& Survival Analysis in Clinical Settings}
\label{sec:survival}

Binary conversion prediction --- \emph{will this MCI subject convert by end of
follow-up?} --- discards information a clinician needs, namely \emph{when}. Survival
analysis models the hazard directly and handles right-censoring principledly: a subject
who has not converted by their last visit is not a negative example, they are an
observation censored at that time. The classical toolkit spans the semi-parametric Cox
proportional hazards model, non-parametric Kaplan--Meier estimation, random survival
forests, and neural extensions such as DeepSurv and DeepHit.

This thesis treats survival as scoped future work rather than a delivered result; the
reasons and the current state of the PROGNOSER component are stated honestly in
Section~\ref{sec:prognoser-scope}.

%==============================================================================
\section{What Is Actually Unsettled: Grey Areas and Research Gaps}
\label{sec:gaps}

The related-work sections above present the field as it presents itself. This section
presents it as the evidence supports. Eight questions are genuinely open; this thesis
takes a position on each, and the position is stated here so that the results chapters can
be read as answers rather than as isolated numbers.

%------------------------------------------------------------------------------
\subsection{Does graph topology carry signal beyond the connectivity matrix itself?}
\label{sec:gap-topology}

The justification for applying GNNs to connectomes is that brain networks have meaningful
topology and that message passing exploits it. The evidence for the second half of that
claim is weak. Lara-Rangel and Heinbaugh evaluated graph transformers on the NeuroGraph
benchmark and found not only that transformers offered no advantage over conventional
GNNs, but --- far more damagingly --- that \emph{both model families retained their
accuracy when every edge was removed}. If predictions survive the deletion of the entire
graph structure, the graph structure was not what the model was using.

This work's pipeline is exposed to the same critique in a specific, checkable way. Edges
are constructed by $k$-nearest-neighbour sparsification with $k = 8$ on $|\mathrm{FC}|$,
and the resulting \texttt{edge\_attr} is a \emph{binary} 0/1 mask: once a correlation has
decided whether an edge exists, its magnitude is discarded. The 200-dimensional node
feature vector, meanwhile, retains the full continuous correlation profile. The graph is
thus a lossy, binarised projection of information the node features already carry in full.

\grey{Is connectome topology informative, or a lossy re-encoding of the node features?}{%
Node features and edges in the standard FC-graph formulation are two views of one matrix,
not two information sources. Under $k$NN-8 binarisation the edge view is strictly poorer
than the node view. Whether message passing over that view can add anything is an open
question the field has largely assumed away; the encoder ablation in
Chapter~\ref{ch:results} is a direct test on this cohort.}

%------------------------------------------------------------------------------
\subsection{Does message aggregation help at all?}
\label{sec:gap-aggregation}

The strongest recent evidence says no. Re-examining graph deep learning against classical
machine learning across four large-scale neuroimaging studies, the authors of
\emph{Rethinking Functional Brain Connectome Analysis} report that the message aggregation
mechanism ``does not help with predictive performance as typically assumed, but rather
consistently degrades it,'' and propose a hybrid linear/graph-attention model in response.
Their closing recommendation --- caution in adopting complex models, and rigorous designs
to establish tangible gains --- is the methodological stance this thesis adopts.

That study is, however, \textbf{cross-sectional and large-$N$}. It says nothing about the
regime that defines this work: longitudinal sequences, $N < 150$, irregular follow-up. The
gap is real and this thesis occupies it.

\gap{The message-aggregation critique has never been tested in the longitudinal, small-$N$
clinical regime.}{%
The published negative result covers four large cross-sectional studies. Whether it
survives when the task is trajectory-based, the cohort is one clinical site, and $N < 150$
is unknown. Chapter~\ref{ch:results}'s four-arm encoder ablation is designed to answer
exactly this, and the ADNI replication in Section~\ref{sec:external-validation} tests
whether the answer is cohort-specific.}

%------------------------------------------------------------------------------
\subsection{What counts as a fair floor at $N < 150$?}
\label{sec:gap-floors}

Reported accuracies for MCI-to-AD conversion prediction from rs-fMRI span roughly
0.85--0.97 in the literature, with several studies claiming combined-modality accuracies
of 97\%. Those numbers are not comparable to each other and frequently not credible.
Connectome-based studies are known to be inflated by leakage through feature selection
performed outside the cross-validation loop, covariate correction fitted on the full
sample, and --- most severely --- non-independence between train and test subjects. Slice-
or scan-level cross-validation, in which multiple scans from the same subject land on both
sides of the split, has been shown to inflate accuracy by 30--55 percentage points.

The deeper problem is that a headline AUC is uninterpretable without a floor. A model
using only age, sex, and visit timing frequently performs respectably on conversion tasks,
because those variables genuinely carry prognostic information. A paper that reports 0.90
without reporting what age and sex alone achieve has not established that its imaging
pipeline contributes anything.

\edge{Every results table in this thesis is anchored against explicit floors.}{%
Two reference floors are computed on the identical splits and reported alongside every
model: a metadata-only baseline (age, sex, visit time --- experiment
\texttt{sanity-metadata-baseline}) and a raw-FC baseline that bypasses the graph encoder
entirely. An arm that does not clear the metadata floor is not evidence about encoders; it
is evidence about the cohort. This is standard practice in principle and rare in the
published brain-GNN literature in fact.}

%------------------------------------------------------------------------------
\subsection{Nobody reports same-seed replicate variance}
\label{sec:gap-determinism}

Deep learning training is non-deterministic at two levels. Software randomness (PRNG seeds,
data-loader shuffling, dropout masks) is what seeding addresses and what the field reports.
\emph{Hardware} non-determinism --- cuDNN algorithm selection and, decisively,
floating-point reduction order in GPU atomics --- is not removed by seeding, and is the
larger term: analyses of ResNet training attribute roughly 80\% of accuracy standard
deviation to non-deterministic sources rather than to seed choice.

PyTorch exposes the remedy: \texttt{torch.use\_deterministic\_algorithms(True)} together
with the \texttt{CUBLAS\_WORKSPACE\_CONFIG} environment variable, set before the CUDA
context is created. Scatter-type backward operations --- which includes the
\texttt{TopKPooling} and index-gather operations at the heart of several brain-graph
architectures --- are precisely the operations that remain non-deterministic until those
flags are set.

The consequence for the literature is stark and, as far as this review found,
unaddressed. A paper reporting ``mean AUC $\pm$ SD over 4 seeds'' has measured
seed-to-seed variance \emph{plus} run-to-run variance, and has attributed the whole of it
to seeds. If the run-to-run term dominates, the reported error bar is not an error bar for
the quantity it names, and any significance test built on it is invalid. Verifying this
requires re-running a \emph{fixed} seed several times --- an experiment that costs nothing
conceptually and that no brain-GNN paper surveyed for this thesis reports.

\gap{Same-seed replicate variance is not reported anywhere in the brain-graph
literature.}{%
Multi-seed spread is universally reported; within-seed spread essentially never is.
Without the latter, the former cannot be interpreted, because the two are confounded.
Chapter~\ref{ch:sota-repair} performs this decomposition for a published MICCAI baseline
across 19 runs and finds within-seed variance \emph{exceeds} between-seed variance ---
which invalidates the error bars, not the direction, of the comparison built on it.}

\edge{The pipeline developed here is bit-reproducible; the audited baseline is not.}{%
The \texttt{none}-arm GELSTM configuration re-runs byte-identically at a fixed seed, so
the non-determinism documented in Chapter~\ref{ch:sota-repair} is a property of the
audited baseline's operations rather than a shared property of the harness. That contrast
is a stronger and more defensible finding than the accuracy margin it sits beside.}

%------------------------------------------------------------------------------
\subsection{Longitudinal evaluation protocols are not standardised --- and the choice
  changes the conclusion}
\label{sec:gap-protocol}

``Early detection'' results are typically presented as discriminative performance versus
the number of visits available, $N = 1, 2, 3, \ldots$. There are two ways to build that
curve and they disagree.

The \emph{variable-cohort} construction keeps, at each $N$, all subjects with at least $N$
visits, truncated to their first $N$. The cohort therefore shrinks as $N$ grows, and each
point on the curve is computed on a different, self-selected population. The
\emph{fixed-cohort} construction restricts to the subjects who have the maximum $N$
throughout, so every point describes the same people with progressively more information.

The difference is not academic. In this thesis's own diagnostics, the static graph
classifier appears to \emph{lose} accuracy from $N = 3$ to $N = 4$ under the
variable-cohort construction ($0.950 \rightarrow 0.914$), and rises monotonically under the
fixed-cohort construction ($0.886 \rightarrow 0.914$). The apparent degradation was the
changing population, not a fourth visit hurting the model. A paper reporting only the
variable-cohort curve would have drawn --- and defended --- a false conclusion.

\gap{No convention exists for constructing AUC-versus-visit-count curves, and the two
common constructions can yield opposite conclusions on identical runs.}{%
This thesis reports both constructions side by side and treats the fixed-cohort version as
the claim-bearing one, with the variable-cohort version shown only to expose the
composition shift. The per-$N$ converter/non-converter counts are printed beside every
curve so the shift is visible rather than inferred.}

%------------------------------------------------------------------------------
\subsection{Follow-up duration is a confound, and it is almost never checked}
\label{sec:gap-attrition}

In any real longitudinal cohort, how long a subject is followed is not independent of
their clinical trajectory. In the DELCODE cohort used here, converters are followed
\emph{longer} than stable-MCI subjects: mean 3.43 versus 2.74 scans, Mann--Whitney
$U$ test $p = 0.002$. Sequence length is therefore correlated with the label, and a
longitudinal model has an available shortcut that has nothing to do with brain function.

Detecting this requires care, because the naive check fails. The between-subject
correlation between predicted probability and visit count is approximately zero in every
model examined here --- which looks reassuring and is not. It is a Simpson's paradox: the
marginal trend (converters followed longer) and the within-non-converter trend cancel.

The decisive test is the \emph{within-subject} slope: as one subject's own visits
accumulate, does the model's predicted probability move? A length shortcut would be flat
within-subject. Instead, converters drift upward and 79--86\% of non-converters drift
downward, which is the signature of evidence accumulation rather than a count cue.

\edge{The follow-up-duration confound is measured, decomposed, and refuted, not assumed
away.}{%
The within-subject slope analysis distinguishes evidence accumulation from a
sequence-length shortcut, and resolves a Simpson's paradox that makes the naive
between-subject check misleading. Informative censoring of this kind is near-universally
unreported in conversion-prediction papers despite being structurally present in every
clinical cohort.}

%------------------------------------------------------------------------------
\subsection{Is reconstruction the right pretext task for a connectome?}
\label{sec:gap-pretext}

Self-supervised pretraining works when the pretext task requires the representation to
encode what the downstream task needs. Reconstruction requires the latent to retain
whatever is necessary to rebuild edges. Conversion prediction requires it to retain
whatever distinguishes a converter from a stable subject. These are different
requirements, and nothing guarantees the first implies the second. Between-subject
variance that is diagnostically informative but reconstruction-irrelevant has no reason to
survive a bottleneck optimised for reconstruction.

The alternative objectives now dominant at scale --- contrastive and joint-embedding
predictive objectives such as those in Brain-JEPA --- are motivated in part by exactly
this concern: predicting in latent space rather than reconstructing in input space avoids
spending capacity on high-variance, low-information detail. Whether that reasoning
transfers to $N < 150$ single-site pretraining is untested.

\grey{Whether reconstruction pretraining transfers to clinical discrimination is assumed,
not demonstrated.}{%
The reconstruction objective is inherited across the brain-graph literature largely
unexamined. This thesis tests it with a four-arm ablation that separates the value of the
pretraining, the value of the architecture at random initialisation, and the value of
having an encoder at all --- with the interpretation table pre-registered before any
numbers existed, so the conclusion cannot be fitted to the result.}

%------------------------------------------------------------------------------
\subsection{Cross-cohort visit-time semantics: the silent-drop failure mode}
\label{sec:gap-cohorts}

External validation is the accepted remedy for single-cohort overfitting, and it introduces
a data-representation problem that is easy to get wrong silently. DELCODE encodes visits as
\emph{nominal protocol months} drawn from a small discrete set ($0, 12, 24, 36, 48, 60$) ---
the scheduled visit label. ADNI and OASIS-3 encode \emph{actual elapsed days from baseline}
(\texttt{ses-d0381}) --- continuous and irregular.

These are different quantities, and conflating them fails in both directions. Collapsing
elapsed days onto the nearest protocol month ($381/30.44 \approx 13$) discards precisely
the irregularity that a $\Delta t$-conditioned model exists to exploit; keeping days
continuous breaks any allow-list filter built on protocol-month set membership. Worse, the
naive failure is silent: a month parser written for DELCODE filenames returns
\texttt{None} on every ADNI and OASIS-3 session, each subject lands with zero visits, and
the cohort disappears without raising anything. The result is a clean, wrong, empty
evaluation.

\gap{Harmonising visit-time semantics across cohorts is treated as preprocessing
bookkeeping, but it is the canonical train-uniform / test-irregular fairness failure.}{%
Training on regularly spaced protocol months and evaluating on irregular elapsed times is
a distribution shift in the model's own input, not a nuisance. This thesis separates the
three concepts explicitly --- \emph{visit index} for ordering, \emph{protocol month} for
leakage filtering, \emph{elapsed months} for model input --- with one function per concept
dispatching on cohort, and never guesses a protocol month it cannot confidently assign.}

%------------------------------------------------------------------------------
\subsection{Summary of positions}

\begin{table}[htbp]
  \centering
  \small
  \begin{tabular}{@{}p{0.29\linewidth}p{0.31\linewidth}p{0.31\linewidth}@{}}
    \toprule
    \textbf{Open question} & \textbf{Field's default assumption} & \textbf{Where this thesis lands} \\
    \midrule
    Topology beyond node features (\S\ref{sec:gap-topology})
      & Graph structure is informative
      & Tested directly; $k$NN-8 binarisation is lossy by construction \\
    Message aggregation (\S\ref{sec:gap-aggregation})
      & Aggregation helps
      & Published large-$N$ evidence says it degrades; tested here longitudinally \\
    Evaluation floors (\S\ref{sec:gap-floors})
      & Headline AUC suffices
      & Metadata and raw-FC floors in every table \\
    Same-seed variance (\S\ref{sec:gap-determinism})
      & Multi-seed SD is the error bar
      & Decomposed over 19 runs; within-seed dominates \\
    Visit-horizon protocol (\S\ref{sec:gap-protocol})
      & Either construction is fine
      & Both reported; fixed-cohort is claim-bearing \\
    Follow-up duration (\S\ref{sec:gap-attrition})
      & Not checked
      & Within-subject slopes; Simpson's paradox resolved \\
    Pretext objective (\S\ref{sec:gap-pretext})
      & Reconstruction transfers
      & Pre-registered four-arm ablation \\
    Cross-cohort visit time (\S\ref{sec:gap-cohorts})
      & Preprocessing detail
      & Three separated concepts, cohort-dispatched \\
    \bottomrule
  \end{tabular}
  \caption{The eight open questions of Section~\ref{sec:gaps}, the assumption the
    literature usually makes about each, and the position this thesis takes.}
  \label{tab:gap-summary}
\end{table}
